% !TEX root = ../../ctfp-print.tex

\lettrine[lhang=0.17]{こ}{れまで}関数型の意味をあいまいにしてきた。関数型は他の型とは異なる。

たとえば\code{Integer}を考えよう：これは単なる整数の集合だ。
\code{Bool}は2要素の集合だ。しかし関数型
$a\to b$はそれ以上のものだ：対象$a$と$b$の間の射の集合だ。
任意の圏における2つの対象の間の射の集合を\newterm{ホム集合}と呼ぶ。
$\Set$圏ではすべてのホム集合はそれ自体が同じ圏における対象となっている――
なぜならそれは結局のところ\emph{集合}だからだ。

\begin{figure}[H]
  \centering
  \includegraphics[width=0.35\textwidth]{images/set-hom-set.jpg}
  \caption{Setにおけるホム集合はただの集合}
\end{figure}

\noindent
他の圏では同じことは言えず、そこではホム集合は圏の外部にある。
それらは\emph{外部}ホム集合と呼ばれる。

\begin{figure}[H]
  \centering
  \includegraphics[width=0.35\textwidth]{images/hom-set.jpg}
  \caption{圏Cにおけるホム集合は外部集合}
\end{figure}

\noindent
関数型が特別なのは、$\Set$圏の自己参照的な性質によるものだ。しかし、少なくとも一部の圏では、
ホム集合を表す対象を構成する方法がある。そのような対象を\newterm{内部}ホム集合と呼ぶ。

\section{普遍的構成}

関数型が集合であることをしばらく忘れて、関数型――より一般には内部ホム集合――をゼロから構成してみよう。
いつものように$\Set$圏をヒントにするが、集合の性質は使わないよう注意する。
そうすることで、この構成が他の圏でも自動的に機能するようになる。

関数型は、引数の型と結果の型との関係から、複合型と見なすことができる。
複合型の構成――対象間の関係に関わるもの――はすでに見てきた。
\hyperref[products-and-coproducts]{積型と余積型}を定義するために普遍的構成を使った。
同じ方法で関数型を定義できる。構成する関数型、引数の型、結果の型という3つの対象が必要だ。

これら3つの型を結びつける明白なパターンは\newterm{関数適用}または\newterm{評価}と呼ばれる。
関数型の候補を$z$と呼ぼう（$\Set$圏でなければ、これは他の対象と同様にただの対象だ）。
引数の型$a$（対象）があれば、適用はこのペアを結果の型$b$（対象）に写す。
3つの対象があり、そのうちの2つ（引数の型と結果の型を表すもの）は固定されている。

また適用という写像もある。この写像をパターンにどう組み込むか？
対象の内部を見ることが許されるなら、関数$f$（$z$の元）と引数$x$（$a$の元）をペアにして、
$f x$（$x$への$f$の適用で$b$の元）に写すことができる。

\begin{figure}[H]
  \centering\includegraphics[width=0.35\textwidth]{images/functionset.jpg}
  \caption{Setでは、関数の集合$z$から関数$f$を選び、集合（型）$a$から引数$x$を選べる。集合（型）$b$の元$f x$が得られる。}
\end{figure}

\noindent
しかし、個々のペア$(f, x)$を扱う代わりに、関数型$z$と引数の型$a$の\emph{積}全体について
話すこともできる。積$z\times{}a$は対象であり、その対象から$b$への矢印$g$を適用射として選べる。
$\Set$では、$g$はすべてのペア$(f, x)$を$f x$に写す関数となる。

つまりパターンはこうだ：2つの対象$z$と$a$の積が射$g$によって別の対象$b$に結ばれている。

\begin{figure}[H]
  \centering
  \includegraphics[width=0.4\textwidth]{images/functionpattern.jpg}
  \caption{普遍的構成の出発点となる対象と射のパターン}
\end{figure}

\noindent
このパターンは普遍的構成を使って関数型を特定するのに十分か？すべての圏でそうというわけではない。
しかし、我々が興味を持つ圏ではそうだ。もう一つの問い：積を定義せずに関数オブジェクトを
定義できるか？積が存在しない圏や、すべての対象ペアに対して積が存在しない圏もある。
答えはノーだ：積型がなければ関数型もない。これについては指数の話をするときに再び触れる。

普遍的構成を復習しよう。対象と射のパターンから始める。これは不正確なクエリであり、
通常は非常に多くのヒットをもたらす。特に$\Set$では、ほぼすべてのものがすべてのものに
つながっている。任意の対象$z$を取り、$a$との積を作れば、$b$への関数が存在する
（$b$が空集合でない限り）。

そこで秘密兵器：ランキングを適用する。これは通常、候補オブジェクト間に一意な写像を要求することで
行われる――その写像が何らかの形で構成を因数分解する。今の場合、$z\times a$から$b$への射$g$を
持つ$z$が、別の$z'$とその適用$g'$より\emph{優れている}のは、$g'$の適用が$g$の適用を
通じて因数分解されるような一意な写像$h$が$z'$から$z$に存在する場合に限る、と定める。
（ヒント：図を見ながらこの文を読もう。）

\begin{figure}[H]
  \centering
  \includegraphics[width=0.4\textwidth]{images/functionranking.jpg}
  \caption{関数オブジェクトの候補の間でのランキングの確立}
\end{figure}

\noindent
ここからが難しいところだ。この普遍的構成をここまで先延ばしにした主な理由でもある。
射$h \Colon z'\to z$が与えられたとき、$z'$と$z$の両方が$a$と掛け合わされた図式を
閉じたい。$z'$から$z$への写像$h$が与えられたとき、本当に必要なのは
$z' \times a$から$z \times a$への写像だ。
\hyperref[functoriality]{積の関手性}について議論した後、その方法がわかった。
積自体が関手（より正確には自己双関手）なので、射のペアを持ち上げることができる。
つまり、対象の積だけでなく射の積も定義できる。

積$z' \times a$の第2成分には触れないので、射のペア$(h, \id)$（$\id$は$a$の恒等射）を
持ち上げることになる。

つまり、ある適用$g'$を別の適用$g$で因数分解する方法はこうだ：
\[g' = g \circ (h \times \id)\]
ここでのキーは射に対する積の作用だ。

普遍的構成の第3の部分は、普遍的に最良の対象を選ぶことだ。この対象を$a \Rightarrow b$と呼ぼう
（これはHaskellの型クラス制約と混同しないよう、一つの対象の記号名と考えてほしい――
後で異なる名前の付け方を議論する）。この対象には固有の適用が付属している――
$(a \Rightarrow b) \times a$から$b$への射で、$\mathit{eval}$と呼ぶ。
対象\code{$a \Rightarrow b$}は、関数オブジェクトの他のすべての候補が、その適用射$g$が
$\mathit{eval}$を通じて因数分解されるような形で一意にそれに写せるとき最良だ。
この対象は我々のランキングによれば他のどの対象よりも優れている。

\begin{figure}[H]
  \centering
  \includegraphics[width=0.4\textwidth]{images/universalfunctionobject.jpg}
  \caption{普遍的関数オブジェクトの定義。上図と同じ図だが、今は対象$a \Rightarrow b$が\emph{普遍的}だ。}
\end{figure}

\noindent
形式的には：

\begin{longtable}[]{@{}l@{}}
  \toprule
  \begin{minipage}[t]{0.97\columnwidth}\raggedright\strut
    $a$から$b$への\emph{関数オブジェクト}とは、射
    \[\mathit{eval} \Colon ((a \Rightarrow b) \times a) \to b\]
    を伴う対象$a \Rightarrow b$であり、任意の他の対象$z$と射
    \[g \Colon z \times a \to b\]
    に対して、$g$を$\mathit{eval}$を通じて因数分解する一意な射
    \[h \Colon z \to (a \Rightarrow b)\]
    が存在するものだ：
    \[g = \mathit{eval} \circ (h \times \id)\]
  \end{minipage}\tabularnewline
  \bottomrule
\end{longtable}

\noindent
もちろん、ある圏において対象の任意のペア$a$と$b$に対してそのような対象$a \Rightarrow b$が
存在する保証はない。しかし$\Set$では常に存在する。さらに$\Set$では、この対象は
ホム集合$\Set(a, b)$と同型だ。

このため、Haskellでは関数型\code{a -> b}を圏論的な関数オブジェクト$a \Rightarrow b$として解釈する。

\section{カリー化}

関数オブジェクトの候補をもう一度見てみよう。しかし今度は、射$g$を$z$と$a$の2変数の
関数として考えよう。
\[g \Colon z \times a \to b\]
積からの射であることは、2変数の関数であることに最も近い。特に$\Set$では、$g$は
集合$z$からの値と集合$a$からの値のペアを引数とする関数だ。

一方、普遍性により、そのような$g$それぞれに対して、$z$を関数オブジェクト$a \Rightarrow b$に
写す一意な射$h$が存在する。
\[h \Colon z \to (a \Rightarrow b)\]
$\Set$では、これは単に$h$が型$z$の変数を1つ取って$a$から$b$への関数を返す関数であることを
意味する。それにより$h$は高階関数となる。したがって普遍的構成は、2変数の関数と
関数を返す1変数の関数の間の一対一対応を確立する。この対応を\newterm{カリー化}と呼び、
$h$を$g$のカリー化バージョンと呼ぶ。

この対応は一対一だ。任意の$g$に対して一意な$h$があり、任意の$h$から常に
次の式で2引数の関数$g$を再構成できる：
\[g = \mathit{eval} \circ (h \times \id)\]
関数$g$は$h$の\emph{非カリー化}バージョンと呼べる。

カリー化はHaskellの構文に本質的に組み込まれている。関数を返す関数：

\src{snippet01}
は2変数の関数として考えられることが多い。括弧を省いたシグネチャをこのように読む：

\src{snippet02}
この解釈は多引数関数を定義する方法にも現れている。例えば：

\src{snippet03}
同じ関数は、関数を返す1引数の関数――ラムダ――として書ける：

\src{snippet04}
これら2つの定義は等価であり、どちらも1つの引数だけに部分適用して1引数の関数を得られる：

\src{snippet05}
厳密に言えば、2変数の関数はペア（積型）を取るものだ：

\src{snippet06}
2つの表現の間の変換は自明で、それを行う2つの（高階）関数は\code{curry}と
\code{uncurry}と呼ばれる：

\src{snippet07}
そして

\src{snippet08}
\code{curry}は関数オブジェクトの普遍的構成における\emph{因数分解関数}であることに注意しよう。
これはこの形式で書き直すと特に明らかだ：

\src{snippet09}
（念のため：因数分解関数は候補から因数分解する関数を生成する。）

C++のような非関数型言語では、カリー化は可能だが自明ではない。C++の多引数関数は、
Haskellでタプルを取る関数に対応すると考えることができる（ただし、話をさらに複雑にすることに、
C++では明示的な\code{std::tuple}を取る関数も、可変引数関数も、初期化子リストを取る関数も
定義できる）。

テンプレート\code{std::bind}を使ってC++の関数を部分適用できる。たとえば、2つの文字列を
引数とする関数があるとする：

\begin{snip}{cpp}
std::string catstr(std::string s1, std::string s2) {
    return s1 + s2;
}
\end{snip}
1つの文字列を引数とする関数を定義できる：

\begin{snip}{cpp}
using namespace std::placeholders;

auto greet = std::bind(catstr, "Hello ", _1);
std::cout << greet("Haskell Curry");
\end{snip}
C++やJavaより関数型に近いScalaは、その中間にある。部分適用されることを想定した関数を
定義するときは、複数の引数リストで定義する：

\begin{snip}{cpp}
def catstr(s1: String)(s2: String) = s1 + s2
\end{snip}
もちろんこれはライブラリ作成者側の相当な先見性や予知力を必要とする。

\section{指数}

数学の文献では、関数オブジェクト、すなわち2つの対象$a$と$b$の間の内部ホムオブジェクトを
\newterm{指数}と呼ぶことが多く、$b^{a}$と表記する。引数の型が指数の位置にあることに注意しよう。
この表記法は最初は奇妙に見えるかもしれないが、関数と積の関係を考えれば完全に理にかなっている。
内部ホムオブジェクトの普遍的構成では積を使わなければならないことはすでに見たが、
つながりはそれよりも深い。

これは有限型――\code{Bool}、\code{Char}、あるいは\code{Int}や\code{Double}のように
有限個の値を持つ型――の間の関数を考えると最もよくわかる。
そのような関数は、少なくとも原理的には、完全にメモ化してデータ構造として参照できる。
そしてこれが、射である関数と対象である関数型の等価性の本質だ。

例えば\code{Bool}からの（純粋な）関数は、値のペア――\code{False}に対応する1つと
\code{True}に対応する1つ――によって完全に特定される。
\code{Bool}から\code{Int}へのすべての可能な関数の集合は\code{Int}のペアの集合だ。
これは積\code{Int} \times{} \code{Int}と同じで、表記法を少し工夫すれば\code{Int}\textsuperscript{2}だ。

別の例として、256個の値を含むC++の\code{char}型を見てみよう（Haskellの\code{Char}は
Unicodeを使うためより大きい）。C++標準ライブラリには、通常ルックアップを使って実装される
いくつかの関数がある。\code{isupper}や\code{isspace}のような関数は、256個のブール値の
タプルと等価なテーブルを使って実装される。タプルは積型なので、256個のブールの積を
扱っている：\code{bool \times{} bool \times{} bool \times{} ... \times{} bool}。
算術から、反復された積はべき乗を定義することがわかっている。
\code{bool}を256回（または\code{char}回）「掛け算」すると、\code{char}乗の\code{bool}、
つまり\code{bool}\textsuperscript{\code{char}}が得られる。

256タプルの\code{bool}として定義された型には何個の値があるか？ちょうど$2^{256}$だ。
これはまた\code{char}から\code{bool}への異なる関数の数でもあり、
各関数は一意な256タプルに対応する。同様に\code{bool}から\code{char}への関数の数は
$256^{2}$と計算できる。関数型の指数表記はこれらの場合に完全に意味をなす。

\code{int}や\code{double}からの関数を完全にメモ化しようとは思わないだろう。
しかし関数とデータ型の等価性は、常に実用的でなくても存在する。
また、リスト、文字列、木などの無限型もある。それらの型からの関数を積極的にメモ化すると
無限のストレージが必要になる。しかしHaskellは遅延評価言語なので、
遅延評価された（無限の）データ構造と関数の境界はあいまいだ。
この関数対データの二重性により、HaskellのMonoidの関数型が圏論的な指数オブジェクトと
同一視されるが、それは我々の\emph{データ}という考えにより近い。

\section{デカルト閉圏}

型と関数のモデルとして集合の圏を使い続けるが、そのために使える圏のより大きなファミリーがある
ことは言及する価値がある。これらの圏は\newterm{デカルト閉圏}と呼ばれ、
$\Set$はそのような圏の一例に過ぎない。

デカルト閉圏は以下を含まなければならない：

\begin{enumerate}
  \tightlist
  \item
        終端対象、
  \item
        任意の対象ペアの積、そして
  \item
        任意の対象ペアの指数。
\end{enumerate}
指数を反復された積（おそらく無限回の）と見なせば、デカルト閉圏を任意のアリティの積を
サポートする圏と考えることができる。特に、終端対象は0個の対象の積――
つまりある対象の0乗――と見なすことができる。

コンピュータサイエンスの観点からデカルト閉圏が興味深いのは、すべての型付きプログラミング言語の
基礎をなす単純型付きラムダ計算のモデルを提供するからだ。

終端対象と積にはそれぞれ双対がある：始対象と余積だ。それら2つもサポートし、積が余積に
分配される：
\begin{gather*}
  a \times (b + c) = a \times b + a \times c \\
  (b + c) \times a = b \times a + c \times a
\end{gather*}
デカルト閉圏を\newterm{双デカルト閉圏}と呼ぶ。次の節では、$\Set$がその代表例である
双デカルト閉圏がいくつかの興味深い性質を持つことを見る。

\section{指数と代数的データ型}

関数型を指数として解釈することは、代数的データ型の枠組みに非常によく合致する。
0と1、和、積、指数に関する高校代数のすべての基本的な恒等式が、双デカルト閉圏では
それぞれ始対象と終対象、余積、積、指数についてほぼそのまま成り立つことがわかる。
これらを証明するツール（随伴やYoneda補題など）はまだないが、貴重な直感の源として
ここに列挙しておく。

\subsection{0乗}

\[a^{0} = 1\]
圏論的解釈では、0を始対象に、1を終対象に、等式を同型に置き換える。
指数は内部ホムオブジェクトだ。この特定の指数は始対象から任意の対象$a$への射の集合を表す。
始対象の定義により、そのような射はちょうど1つ存在するので、ホム集合$\cat{C}(0, a)$は
単元集合だ。単元集合は$\Set$の終端対象なので、この恒等式は$\Set$で自明に成り立つ。
この恒等式が任意の双デカルト閉圏で成り立つと言っている。

Haskellでは、0を\code{Void}に、1を単位型\code{()}に、指数を関数型に置き換える。
\code{Void}から任意の型\code{a}への関数の集合が単位型と等価であると主張している――
つまり単元だ。言い換えれば、\code{Void -> a}という関数は1つしかない。
この関数は以前見たことがある：\code{absurd}と呼ばれる。

これには2つの理由から少し注意が必要だ。1つは、Haskellには実際には非inhabited型がない――
すべての型に「終わらない計算の結果」またはボトムが含まれる。2つ目の理由は、
\code{absurd}のすべての実装が等価だということだ。なぜなら、何をしようと誰もそれを
実行できないからだ。\code{absurd}に渡せる値はない。（そして終わらない計算を渡すことができたとしても、
決して返らない！）

\subsection{1の累乗}

\[1^{a} = 1\]
この恒等式は$\Set$で解釈すると、終端対象の定義を再表明する：任意の対象から終端対象への
一意な射が存在する。一般に、$a$から終端対象への内部ホムオブジェクトは終端対象自体と同型だ。

Haskellでは、任意の型\code{a}から単位型への関数は1つだけだ。
この関数は以前見たことがある――\code{unit}と呼ばれる。
\code{const}を\code{()}に部分適用した関数と考えることもできる。

\subsection{1乗}

\[a^{1} = a\]
これは、終端対象からの射が対象\code{a}の「元」を選ぶために使えるという観察の再表明だ。
そのような射の集合はその対象自体と同型だ。$\Set$でもHaskellでも、同型は集合\code{a}の
元と、それらの元を選ぶ関数\code{() -> a}の間にある。

\subsection{和の指数}

\[a^{b+c} = a^{b} \times a^{c}\]
圏論的には、2つの対象の余積からの指数は2つの指数の積と同型と言っている。
Haskellでは、この代数的恒等式は非常に実用的な解釈を持つ。2つの型の和からの関数は
個々の型からの関数のペアと等価であると教えてくれる。これはまさに和に対する関数を定義するときに
使う場合分けだ。\code{case}節を使って1つの関数定義を書く代わりに、通常は各型コンストラクタを
個別に扱う2つ（以上）の関数に分割する。例えば、和型\code{(Either Int Double)}からの関数を取ろう：

\src{snippet10}
それぞれ\code{Int}と\code{Double}からの関数のペアとして定義できる：

\src{snippet11}
ここで\code{n}は\code{Int}、\code{x}は\code{Double}だ。

\subsection{指数の指数}

\[(a^{b})^{c} = a^{b \times c}\]
これは指数オブジェクトの言葉でカリー化を表現しているに過ぎない。
関数を返す関数は、積からの関数（2引数の関数）と等価だ。

\subsection{積の指数}

\[(a \times b)^{c} = a^{c} \times b^{c}\]
Haskellでは：ペアを返す関数は、それぞれペアの一方の要素を生成する関数のペアと等価だ。

これらのシンプルな高校代数の恒等式が圏論に持ち上げられ、関数型プログラミングに実用的な
応用を持つことは驚くべきことだ。

\section{Curry-Howard同型}

論理と代数的データ型の対応についてはすでに触れた。\code{Void}型と単位型\code{()}は
偽と真に対応する。積型と和型は論理的連言$\wedge$（AND）と選言$\vee$（OR）に対応する。
この枠組みでは、定義したばかりの関数型は論理的含意$\Rightarrow$に対応する。
言い換えれば、型\code{a -> b}は「aならばb」と読める。

Curry-Howard同型によれば、すべての型は命題――真か偽かもしれない文や判断――として解釈できる。
そのような命題は型がinhabitedであれば真、そうでなければ偽と見なされる。特に、論理的含意は
対応する関数型がinhabitedであれば真、つまりその型の関数が存在すれば真だ。
関数の実装は定理の証明だ。プログラムを書くことは定理を証明することと等価だ。
いくつかの例を見てみよう。

関数オブジェクトの定義で導入した関数\code{eval}を取り上げよう。そのシグネチャは：

\src{snippet12}
関数とその引数からなるペアを取り、適切な型の結果を生成する。これは射のHaskell実装だ：

\[\mathit{eval} \Colon (a \Rightarrow b) \times a \to b\]
これは関数型$a \Rightarrow b$（または指数オブジェクト$b^{a}$）を定義する。
Curry-Howard同型を使ってこのシグネチャを論理的命題に翻訳しよう：

\[((a \Rightarrow b) \wedge a) \Rightarrow b\]
この文の意味はこうだ：$b$が$a$から導かれることが真であり、$a$が真なら、$b$は真でなければならない。
これは完全に直感的で、古代から\newterm{モーダスポネンス}として知られている。
この定理は関数を実装することで証明できる：

\src{snippet13}
\code{a}を取り\code{b}を返す関数\code{f}と型\code{a}の具体的な値\code{x}のペアを
渡してもらえれば、単純に関数\code{f}を\code{x}に適用することで型\code{b}の具体的な値を
生成できる。この関数を実装することで、型\code{((a -> b), a) -> b}がinhabitedであることを
示した。したがって\newterm{モーダスポネンス}は我々の論理で真だ。

明らかに偽の命題はどうか？たとえば：$a$または$b$が真ならば$a$は真でなければならない。

\[a \vee b \Rightarrow a\]
これは明らかに間違いだ。$a$が偽で$b$が真を選べば、それが反例になる。

Curry-Howard同型を使ってこの命題を関数シグネチャに写すと：

\src{snippet14}
どれだけ試みても、この関数は実装できない――\code{Right}の値で呼び出されたとき、
型\code{a}の値を生成できない。（\emph{純粋}関数の話をしていることを思い出そう。）

最後に、\code{absurd}関数の意味に至る：

\src{snippet15}
\code{Void}が偽に翻訳されることを考えると：

\[\mathit{false} \Rightarrow a\]
偽からは何でも導かれる（\emph{ex falso quodlibet}）。Haskellでのこの文（関数）の
可能な証明（実装）の1つを示す：

\begin{snip}{haskell}
absurd (Void a) = absurd a
\end{snip}
ここで\code{Void}は次のように定義される：

\begin{snip}{haskell}
newtype Void = Void Void
\end{snip}
いつものように、型\code{Void}は扱いにくい。この定義は値を構成することを不可能にする。
なぜなら値を構成するためには一つが必要だからだ。したがって、関数\code{absurd}は
決して呼び出せない。

これらはすべて興味深い例だが、Curry-Howard同型に実用的な側面はあるか？
おそらく日常的なプログラミングにはない。しかし、AgdaやRocq\footnote{以前はCoqとして知られていた}
のようなプログラミング言語があり、定理を証明するためにCurry-Howard同型を活用している。

コンピュータは数学者の仕事を助けているだけでなく、数学の基盤そのものを革命的に変えている。
その分野の最新の研究トピックはホモトピー型理論と呼ばれ、型理論から生まれた。
ブール値、整数、積と余積、関数型などに満ちている。そして疑いを払拭するかのように、
この理論はRocqとAgdaで定式化されている。コンピュータは多くの意味で世界を革命している。

\section{参考文献}

\begin{enumerate}
  \tightlist
  \item
        Ralf Hinze, Daniel W. H. James,
        \urlref{http://www.cs.ox.ac.uk/ralf.hinze/publications/WGP10.pdf}{Reason
          Isomorphically!}. この論文には、本章で触れた圏論における高校代数の恒等式の
        すべての証明が含まれている。
\end{enumerate}
